% Checkpoint35 proposed source; checkpoint34 sources remain unchanged.
\statement{A.113}{Proved}{The entire first-kernel cyclic cohomology is its pure fixed group, and a central sign twist preserves the quotient negative-frame equations}
Retain \(N_1\), \(E_1=F/N_1\), and
\(\rho=\operatorname{Ad}(r)\) in their actual original markings.
There is a canonical pointed homeomorphism
\[
 \partial:T=C_{E_1}(r)=(E_1')^\rho
       \ \xrightarrow{\ \sim\ }\ H^1(C_3,N_1),\qquad
 \partial(h)=[w^{-1}\rho(w)],\quad w\in F',\ wN_1=h.
 \tag{A.113.1}
\]
There is an exact central extension of compact groups
\[
 1\longrightarrow K^\rho\longrightarrow T
   \longrightarrow C_{\mathcal C'}(J)\longrightarrow1,\qquad
 K=(N\cap F')/N_1\ \text{pro-two}.
 \tag{A.113.2}
\]
The natural map \(H^1(C_3,N_1)\to H^1(C_3,K_3')\) is,
under A.109 and (A.113.1), exactly \(q_3:T\to\{I,-I\}\).
Here \(K_3'=\ker(q_3|_{F'})\).

For the ORIGINAL LM7\_PUSH~E.61/E.65 pure \(r\)-frame \(A_r\),
all pure quotient frames of the same action are
\(\overline A_rT\). The class of the exact transporter defect of
a lift of \(\overline A_rh\) is \(\partial(h)\).
Only \(h=1\) admits the original full transporter.

There is moreover an actual central involution
\[
 c_*\in Z(B/N_1)\cap E_1',\qquad
 \alpha(c_*)=c_*,\quad q_4(c_*)=I,\quad q_3(c_*)=-I.
 \tag{A.113.3}
\]
For the specified negative frame \(J_e^\alpha\) in A.111,
both \(J_e^\alpha\) and \(J_e^\alpha c_*\) obey the same two
quotient negative-frame equations and induce the same
\(\kappa_e\) on \(T_e\).
Their coherent classes differ by multiplication by \(c_*\);
this is a fixed-point-free involution of the compact cohomology
space, toggling its three-adic coordinate. It does not change the
original full-group frame, and does not decide the actual CM value.

\begin{proof}
\emph{A.113.1. Both boundary directions and their topology.}
A.109 supplies \((F')^\rho=1\) and \(H^1(C_3,F')=\{*\}\).
We use the convention
\[
 z\rho(z)\rho^2(z)=1,\qquad z\sim n^{-1}z\rho(n).
\]
For \(h\in(E_1')^\rho\), its lift \(w\in F'\) has defect
\(c_w=w^{-1}\rho(w)\in N_1\), whose three-factor norm telescopes.
Replacing \(w\) by \(wn\) gives
\(c_{wn}=n^{-1}c_w\rho(n)\).
Conversely every \(N_1\)-cocycle is an \(F'\)-coboundary by
A.109 and hence comes from such an \(h\).
If \(c_{w_2}=n^{-1}c_{w_1}\rho(n)\), then
\[
                 \rho(w_1nw_2^{-1})=w_1nw_2^{-1}.
\]
The full fixed group is trivial, so \(w_1n=w_2\) and the
two quotient elements are equal. This proves bijectivity.

The compact closed lift set
\[
       L=\{w\in F':w^{-1}\rho(w)\in N_1\}
\]
surjects onto \((E_1')^\rho\). Its continuous defect map followed
by the cocycle orbit map is constant on every lift fiber and
therefore descends continuously. The cocycle space is compact,
and its compact conjugation action has a closed orbit relation;
its quotient is Hausdorff. Thus the bijection is a homeomorphism.
No continuous section of \(L\to(E_1')^\rho\) is assumed.
Finally \(E_1^{\rm ab}=\widehat{\mathbb Z}^2\), and the cyclic
abelian matrix has determinant-three difference from the identity.
Its fixed module is zero, proving \(C_{E_1}(r)=(E_1')^\rho\).

\emph{A.113.2. The exact signed comparison.}
A.106 gives the central pro-two kernel \(K\) in
\[
                 1\to K\to E_1'\to\mathcal C'\to1.
\]
On \(\mathcal C'\) the actual cyclic action is
\(\operatorname{Ad}(J)\); its extra exponent-parity sign vanishes.
For a fixed element downstairs, a lift \(g\) has
\(g^{-1}\rho(g)\) a \(K\)-cocycle.
The complete compact coprime-action lemma in A.109 gives
\(H^1(C_3,K)=1\); writing the defect as \(k^{-1}\rho(k)\)
makes \(gk^{-1}\) fixed. This proves right surjectivity in
(A.113.2). Its kernel is exactly \(K^\rho\).
No splitting of this central extension is asserted.

Use the same \(w\in F'\) in the two boundary constructions for
\(N_1\) and \(K_3'\). The latter has label \(q_3(w)\) by A.109.
Thus the comparison is precisely \(q_3(h)\), without an inverse
or an added sign. It is onto, since A.110 supplies an actual
fixed element with value \(-I\). Its two fibers are the two
closed cosets of \(\ker(q_3|_T)\).
The integral matrix \(J\) and the word \(XYX^{-1}\) are
identified here only after reduction modulo three.

\emph{A.113.3. Every pure quotient frame and its obstruction.}
The original full frame satisfies
\(\rho(A_r)=\phi A_r\), with \(A_r\in F'\).
Two pure quotient \(r\)-frames differ on the right by \(T\).
If \(w\) lifts \(h\in T\), put \(t=A_rw\). Direct multiplication
gives the exact original transporter defect
\[
       t^{-1}\phi^{-1}\rho(t)=w^{-1}\rho(w).
 \tag{A.113.4}
\]
Its class is consequently \(\partial(h)\); it is the basepoint
exactly for \(h=1\). Every element of \(\overline A_rT\) has the
same zero abelianization. This proves that all such lifting
obstructions, rather than only one finite sign, are accounted for.

\emph{A.113.4. The central sign twist in the actual product.}
In (A.112.1)--(A.112.2) define \(c_*\) to have identity dyadic
coordinate, coordinate \(-I\) in \(G_3'\), identity at every
\(p\geq5\), and zero abelian coordinates. The \(\chi\)-value
of \(-I\) is zero, so these are actual allowed coordinates.
It is a nonidentity central involution in \(E_1\).
It also commutes with the original \(s\) and \(r\).
Their actions on the derived odd factors are matrix conjugations,
which fix supported scalars; their extra signs are zero on the
derived group. Their dyadic action fixes its identity coordinate.
Thus it is central in all \(B/N_1\).
The actual action of \(\alpha\) on the odd factors is the matrix
conjugation in A.104, and its dyadic action descends because it
preserves \(M_2\) and \(M_{21}\). Hence \(\alpha(c_*)=c_*\).
This proves (A.113.3), including \(q_4(c_*)=I\).

Let \(J_0\) denote the image of the ORIGINAL \(J_e^\alpha\).
It satisfies
\[
 \alpha(R_e)=J_0R_e^{-1}J_0^{-1},\qquad \alpha(J_0)J_0=1.
\]
Centrality and \(\alpha(c_*)=c_*=c_*^{-1}\) prove the same
two equations for \(J_0c_*\). They preserve the same
\(S_3\)-coset, abelian normalization, and modulo-four frame.
The pure quotient change \(\overline A_r\mapsto\overline A_rc_*\)
effects exactly this change in both original formulas
\(J_e^\alpha=s\phi A_r\) for \(e=1\) and \(J_e^\alpha=A_rs\)
for \(e=-1\).

The induced \(\kappa_e\) is unchanged. For every actual
positive frame \(h\) of A.108,
\[
 \begin{split}
 w_h&=J_0^{-1}\alpha(h)^{-1}vh,\\
 w_h'&=(J_0c_*)^{-1}\alpha(h)^{-1}vh=c_*w_h.
 \end{split}
 \tag{A.113.5}
\]
Multiplication by the central \(\kappa_e\)-fixed involution
preserves cocycles and their equivalences. It induces a continuous
involution of \(H^1(C_2,T_e;\kappa_e)\).
It has no fixed point: the exact \(q_3\) character has trivial
action, no nontrivial coboundaries, and sends the two classes to
opposite signs. By A.112 it toggles only the three-adic binary
coordinate. It is a homeomorphism, not a pointed automorphism.

The alternative frame is not the image of the unique ORIGINAL
full frame. Indeed \(c_*\) has no \(\rho\)-fixed lift in \(F'\),
and \(\partial(c_*)\ne*\) is its exact obstruction.
The finite signed image is the already executed A.110
\(h_{12}=5I\) certificate. Existence of \(c_*\) here follows
from the actual compact product, not from that finite computation.
The original CM parity remains open.
\end{proof}

\begingroup\footnotesize\raggedright
Complete retained proof sources:
\path{work_DE/entire_cyclic_congruence_kernel_cohomology34.md}
and \path{work_DE/central_sign_twist_negative_frames34.md}.
Both complete independent reviews and the A.110 prescribed
finite certificate are preserved. All cohomology spaces here
use actual compact kernels, not abstract completions of
infinite-index discrete subgroups.\par\endgroup
